% !TeX program = xelatex
\documentclass[12pt]{nextbook}

\UseTemplateSet{
  features = {image}
}

\title{English Book Template}
\subtitle{A concise subtitle below the main title}
\author{Author Name}
\date{\today}

\begin{document}

\frontmatter
\maketitle

\chapter*{Preface}

This template is intended for book-length writing, lecture notes, or extended
research manuscripts. The preface can explain motivation, scope, and overall
chapter design.

\tableofcontents

\mainmatter

\chapter{Introduction}

\section{Research Context}

Use this section to introduce the broader context, the existing conversation,
and the reason the book is organized in its present form.

\section{Sources and Method}

This section can describe the corpus, primary materials, editorial choices, or
analytical approach used throughout the manuscript.

\chapter{Sample Chapter}

\section{Long-Form Structure}

Unlike a short article starter, a book template should already assume chapter
progression and a longer reading arc. This file keeps that skeleton in place.

\section{Figure Placeholder}

\begin{figure}[htbp]
  \centering
  \rule{0.7\linewidth}{0.28\linewidth}
  \caption{Chapter figure placeholder}
\end{figure}

\backmatter

\chapter*{Afterword}

This can be replaced with a conclusion, an afterword, or a note on future
development.

\end{document}
